% AI-translated component: printed 272--281.
% Principal PDF 279--288; Google PDF 283--292.
\begingroup
\clearpage
\noindent272\hfill\textsc{Posthumous Papers.}\hfill\mbox{}\par\medskip
\[
\setlength{\arraycolsep}{3pt}
\begin{array}{c@{\quad}rrr@{\quad}rrrrrr@{\;}c@{\;}rrr}
4\,.\,5&31^\circ&36'&57''{,}7&\ast&\ast&-c&+d&&&=&31^\circ&36'&56''{,}275\\
4\,.\,6&102&30&12{,}2&\ast&\ast&-c&\ast&+e&&=&102&30&13{,}050\\
4\,.\,7&165&39&25{,}6&\ast&\ast&-c&\ast&\ast&+f&=&165&39&26{,}450\\[7pt]
5\,.\,6&70&53&14{,}5&\ast&\ast&\ast&-d&+e&&=&70&53&14{,}900\\
5\,.\,7&134&2&27{,}9&\ast&\ast&\ast&-d&\ast&+f&=&134&2&29{,}475\\[7pt]
6\,.\,7&63&9&13{,}4&\ast&\ast&\ast&\ast&-e&+f&=&63&9&13{,}255.
\end{array}
\]
\begin{center}[Error expressions for the observed angles:]\end{center}
\[
\renewcommand{\arraystretch}{1.15}
\begin{array}{rrrrrrr}
+0''{,}625&+a&&&&&\\
+0{,}675&\ast&+b&&&&\\
-0{,}350&\ast&\ast&+c&&&\\
+0{,}425&\ast&\ast&\ast&+d&&\\
+0{,}100&\ast&\ast&\ast&\ast&-e&\\
+0{,}200&\ast&\ast&\ast&\ast&\ast&-f\\[7pt]
+0{,}163&-a&+b&&&&\\
-1{,}000&-a&\ast&+c&&&\\
+0{,}335&-a&\ast&\ast&+d&&\\
\llap{\raisebox{-.6ex}{\text{\tiny.}}\hspace{1.3cm}}+0{,}800&-a&\ast&\ast&\ast&+e&\\
-0{,}940&+a&\ast&\ast&\ast&\ast&-f\\[7pt]
+0{,}075&\ast&-b&+c&&&\\
-1{,}500&\ast&-b&\ast&+d&&\\
-0{,}450&\ast&-b&\ast&\ast&+e&\\
-1{,}750&\ast&+b&\ast&\ast&\ast&-f\\[7pt]
+1{,}425&\ast&\ast&-c&+d&&\\
-0{,}850&\ast&\ast&-c&\ast&+e&\\
-0{,}850&\ast&\ast&-c&\ast&\ast&+f\\[7pt]
-0{,}400&\ast&\ast&\ast&-d&+e&\\
-1{,}575&\ast&\ast&\ast&-d&\ast&+f\\[7pt]
+0{,}145&\ast&\ast&\ast&\ast&-e&+f.
\end{array}
\]

\clearpage
\noindent\mbox{}\hfill\textsc{Station Adjustments.}\hfill273\par\medskip
\begin{center}Normal equations.\end{center}
\[
\begin{array}{r@{\;}c@{\;}rrrrrrr}
0&=&-0{,}613&+6a&-b&-c&-d&-e&-f\\
0\mathord{\cdot}&=&+0{,}963&-a&+6b&-c&-d&-e&-f\\
0&=&-1{,}000&-a&-b&+6c&-d&-e&-f\\
0&=&+2{,}660&-a&-b&-c&+6d&-e&-f\\
0&=&-1{,}145&-a&-b&-c&-d&+6e&-f\\
0&=&+0{,}210&-a&-b&-c&-d&-e&+6f
\end{array}
\]
[from which it follows, since $0=+1{,}075+a+b+c+d+e+f$ results:
\[
\begin{array}{r@{\;}c@{\;}r@{\qquad}r@{\;}c@{\;}r@{\qquad}r@{\;}c@{\;}r}
a&=&-0''{,}066&c&=&-0''{,}011&e&=&+0''{,}010\\
b&=&-0{,}291&d&=&-0{,}534&f&=&-0{,}183,
\end{array}
\]
and hence the mean angular error: $\pm0''{,}96$].

Remarks. The angles were measured with the old 8-inch \textsc{Reichenbach} theodolite, but I sighted with my own telescope of 40-fold magnification; each individual angle was repeated 20 times, with readings from one set of 5 to the next set of 5, the first 5 measured forward and the next 5 backward. The vernier circle was clamped and the telescope was set with the right hand (vernier circle clamp on the right). The movement of the arm with the check telescope was prevented entirely by tightening the screws that pass through the strap; the ball and nut of the micrometer screw on the clamp of the graduated circle were pressed together so firmly that at least no play remained; the ball and nut of the micrometer screw of the vernier circle were clamped not quite so firmly, but firmly enough. The vernier circle was clamped with great care, and only so tightly that the momentum from turning the entire instrument could not shift the vernier circle. The vertical thread and the optical axis were adjusted almost every day; the leveling remained virtually unchanged throughout, and the instrument was mounted very firmly.\}
\vfill
\noindent{\small IX.}\hfill35

\clearpage
\noindent\llap{\raisebox{.2ex}{\tiny.}}274\hfill\textsc{Posthumous Papers.}\hfill\mbox{}\par\medskip
\begin{center}
[5.]\\[3pt]
[Observations on] Breithorn [1822].\\[4pt]
Sept. 13. Afternoon.\\[4pt]
Z[enith-]D[istance] of the surface of the Hauselberg stone.
\end{center}
\noindent[Rep. readings on 4 verniers]
\[
\setlength{\arraycolsep}{3pt}
\GaussRightBrace{\begin{array}{r|rrlcrrr}
0&101^\circ&41'&22''.22''.22''.24''&=&101^\circ&41'&22''{,}50\\
20&102&17&22.23.20.22&=&102&17&21{,}75
\end{array}}\quad90^\circ\ 1'\ 47''{,}9625.
\]
\begin{center}Z.-D. of the surface of the Eschede stone[\textsuperscript{*})].\end{center}
\[
\setlength{\arraycolsep}{3pt}
\GaussRightBrace{\begin{array}{r|rrlcrrr}
0&102&17&22.\ \text{etc.}&=&102&17&21{,}75\\
20&105&49&28.30.30.29&=&105&49&29{,}25
\end{array}}\quad90\ 10\ 36{,}375.
\]
\begin{center}
Theodolite. Circle on the left.\\[4pt]
Eschede, stone --- Hauselberg, stone.
\end{center}
\[
\setlength{\arraycolsep}{3pt}
\GaussRightBrace{\begin{array}{r|rrlcrrr}
0&283&13&63.59.70.65&=&283&14&4{,}25\\
10&62&35&21.24.29.19&=&62&35&23{,}25
\end{array}}\quad121\ 56\ 7{,}900.
\]
\begin{center}
Eschede, stone --- Falkenberg, stone (difficult to distinguish from a nearby treetop).
\end{center}
\[
\setlength{\arraycolsep}{3pt}
\GaussRightBrace{\begin{array}{r|rrlcrrr}
0&62&35&21.\ \text{etc.}&=&62&35&23{,}25\\
10&278&59&39.38.47.43&=&278&59&41{,}75
\end{array}}\quad93\ 38\ 25{,}850.
\]
\begin{center}Eschede, stone --- Hauselberg, stone.\end{center}
\[
\setlength{\arraycolsep}{3pt}
\begin{array}{r|rrlcrrr}
0&278&59&39.\ \text{etc.}&=&278&59&41{,}75\\
10&58&21&15.19.20.14&=&58&21&17{,}00\\
20&197&42&39.37.39.46&=&197&42&40{,}25
\end{array}
\quad
\GaussRightBrace{\begin{array}{rrr}
121&56&9{,}525\\&&8{,}325
\end{array}}\quad121^\circ56'\,8''{,}925.
\]
\begin{center}Wilsede signal tree --- Hauselberg, stone.\end{center}
\[
\setlength{\arraycolsep}{3pt}
\GaussRightBrace{\begin{array}{r|rrlcrrr}
0&197&42&39.\ \text{etc.}&=&197&42&40{,}25\\
2&142&46&40.38.37.38&=&142&46&38{,}25
\end{array}}\quad27\ 28\ 1{,}000.
\]
\vfill\noindent\rule{\textwidth}{.4pt}\par
{\footnotesize [\textsuperscript{*}) The point Eschede is later called Scharnhorst.]}

\clearpage
\noindent\mbox{}\hfill\textsc{Station Adjustments.}\hfill275\par\bigskip
\begin{center}
Sept. 14. Morning, $11^{\mathrm h}$. Very strong air shimmer.\\[4pt]
Z.-D. of the Wilsede heliotrope.
\end{center}
\[
\setlength{\arraycolsep}{3pt}
\begin{array}{r|rrlcrrr}
0&108^\circ&47'&14''.11''.17''.17''&=&108^\circ&47'&14''{,}75\\
2&289&2&45.50.49.47&=&289&2&47{,}75\\
4&109&17&29.35.33.31&=&109&17&32{,}00\\
6&289&33&29.38.35.31&=&289&33&33{,}25
\end{array}
\quad
\GaussRightBrace{\begin{array}{rrr}
90^\circ&7'&46''{,}500\\&7&22{,}125\\&8&0{,}625
\end{array}}\quad90^\circ\ 7'\ 43''{,}083.
\]
\begin{center}Z.-D. of the Eschede heliotrope.\end{center}
\[
\setlength{\arraycolsep}{3pt}
\begin{array}{r|rrlcrrr}
0&289&33&29.\ \text{etc.}&=&289&33&33{,}25\\
4&290&18&29.38.33.30&=&290&18&32{,}50\\
8&291&4&5.\ 9.\ 6.\ 6&=&291&4&6{,}50\\
12&291&48&23.30.24.24&=&291&48&25{,}25
\end{array}
\quad
\GaussRightBrace{\begin{array}{rrr}
90&11&14{,}812\\&11&23{,}500\\&11&4{,}687
\end{array}}\quad90\ 11\ 14{,}333.
\]
\begin{center}
Z.-D. of the Wilsede heliotrope (from $2^{\mathrm h}$ onward). Air somewhat better.
\end{center}
\[
\setlength{\arraycolsep}{3pt}
\begin{array}{r|rrlcrrr}
0&291&48&23.\ \text{etc.}&=&291&48&25{,}25\\
8&292&42&19.23.18.17&=&292&42&19{,}25\\
20&293&58&30.33.34.31&=&293&58&32{,}00
\end{array}
\quad
\GaussRightBrace{\begin{array}{rrr}
90&6&44{,}250\\&6&21{,}062
\end{array}}\quad90\ 6\ 30{,}337.
\]
\begin{center}
Theodolite. Circle on the right.\\[4pt]
Wilsede [heliotrope] --- Hauselberg, stone.
\end{center}
\[
\setlength{\arraycolsep}{3pt}
\GaussRightBrace{\begin{array}{r|rrlcrrr}
0&322&46&31.32.34.35&=&322&46&33{,}00\\
5&185&30&26.22.17.23&=&185&30&22{,}00
\end{array}}\quad27\ 27\ 14{,}200.
\]
\begin{center}Wilsede, heliotrope --- Eschede, stone.\end{center}
\[
\setlength{\arraycolsep}{3pt}
\GaussRightBrace{\begin{array}{r|rrlcrrr}
0&185&30&26.\ \text{etc.}&=&185&30&22{,}00\\
4&307&57&4.\ 6.10.\ 9&=&307&57&7{,}25
\end{array}}\quad149\ 23\ 18{,}6875.
\]
\begin{center}Wilsede, heliotrope --- Hauselberg, stone.\end{center}
\[
\setlength{\arraycolsep}{3pt}
\GaussRightBrace{\begin{array}{r|rrlcrrr}
0&307&57&4.\ \text{etc.}&=&307&57&7{,}25\\
5&170&41&6.\ 9.\ 1.\ 7&=&170&41&5{,}75
\end{array}}\quad27\ 27\ 12{,}300.
\]
\vfill\noindent\hfill35\textsuperscript{$\ast$}\hspace*{2cm}

\clearpage
\noindent276\rlap{\hspace{.12em}\raisebox{.8ex}{\tiny.}}\hfill\textsc{Posthumous Papers.}\hfill\mbox{}\par\medskip
\begin{center}Wilsede [heliotrope] --- Eschede, heliotrope.\end{center}
\[
\setlength{\arraycolsep}{3pt}
\begin{array}{r|rrlcrrr}
0&170^\circ&41'&6''.\ \text{etc.}&=&170^\circ&41'&5''{,}75\\
2&231&54&30.24.26.30&=&231&54&27{,}50\\
12&178&0&68.63.59.67&=&178&1&4{,}25\\
19&212&17&58.53.54.59&=&212&17&56{,}00\\
26&246&34&37.30.35.33&=&246&34&33{,}75\\
33&280&50&59.57.59.58&=&280&50&58{,}25\\
35&342&4&14.20.17.18&=&342&4&17{,}25\\
38&253\rlap{\hspace{.13em}\raisebox{.9ex}{\text{\tiny.}}}&54&20.14.19.14&=&253&54&16{,}75\\
39&104&30&60.56.51.49&=&104&30&54{,}00
\end{array}
\quad
\GaussRightBrace{\begin{array}{rrr}
149^\circ&23'&19''{,}125\\&&20{,}325\\&&18{,}321\\&&20{,}321\\
&&22{,}214\\&&20{,}500\\&&20{,}167\\&&22{,}750
\end{array}}\quad149^\circ23'\,20''{,}300.
\]
\noindent\rule{\textwidth}{.4pt}
\begin{center}
Sept. 15. Theodolite. Circle on the left.\\[4pt]
Wilsede [heliotrope] --- Eschede, stone. Paper strip $0{,}0475\,\mathrm m$ W. of the ctr., amounts to $0''{,}873$.
\end{center}
\[
\setlength{\arraycolsep}{3pt}
\GaussRightBrace{\begin{array}{r|rrlcrrr}
0&284&30&48.51.61.54&=&284&30&53{,}50\\
10&230&37&58.52.59.59&=&230&37&57{,}00
\end{array}}\quad149\ 23\ 17{,}650.
\]
\begin{center}Eschede (as above) --- Falkenberg, stone.\end{center}
\[
\setlength{\arraycolsep}{3pt}
\GaussRightBrace{\begin{array}{r|rrlcrrr}
0&230&37&58.\ \text{etc.}&=&230&37&57{,}00\\
8&259&45&13.13.22.17&=&259&45&16{,}25
\end{array}}\quad93\ 38\ 24{,}906.
\]
\begin{center}Eschede (as above) --- Hauselberg, stone.\end{center}
\[
\setlength{\arraycolsep}{3pt}
\GaussRightBrace{\begin{array}{r|rrlcrrr}
0&259&45&13.\ \text{etc.}&=&259&45&16{,}25\\
5&149&25&49.51.47.50&=&149&25&49{,}25
\end{array}}\quad121\ 56\ 6{,}600.
\]
\noindent\rule{\textwidth}{.4pt}
\begin{center}
Sept. 16.\\[4pt]
Z.-D. of the Falkenberg heliotrope. The first 20 [rep.] in the morning from $9\frac34$---$10\frac34{}^{\mathrm h}$, the rest resumed in the afternoon at $2^{\mathrm h}40^{\mathrm m}$.
\end{center}
\[
\setlength{\arraycolsep}{3pt}
\begin{array}{r|rrlcrrr}
0&293&58&29.36.34.29&=&293&58&32{,}00\\
2&114&3&49.54.50.51&=&114&3&51{,}00\\
4&294&9&19.23.19.18&=&294&9&19{,}75\\
6&114&14&50.50.48.48&=&114&14&49{,}00\\
8&294&20&26.33.27.25&=&294&20&27{,}75\\
10&114&25&54.58.58.57&=&114&25&56{,}75\\
12&294&31&35.35.35.36&=&294&31&35{,}25\\
20&294&53&50.58.55.51&=&294&53&53{,}50\\
22&114&59&14.18.14.14&=&114&59&15{,}00\\
30&115&20&41.50.42.43&=&115&20&44{,}00
\end{array}
\quad
\GaussRightBrace{\begin{array}{rrr}
90&2&39{,}500\\&&44{,}375\\&&44{,}625\\&&49{,}375\\&&44{,}500\\
&&49{,}250\\&&47{,}281\\&&40{,}750\\&&41{,}125
\end{array}}\quad90\ 2\ 44{,}400.
\]

\clearpage
\noindent\mbox{}\hfill\textsc{Station Adjustments.}\hfill277\par\medskip
\begin{center}
Theodolite. Circle on the right.\\[4pt]
Falkenberg, heliotrope --- Wilsede, heliotrope.
\end{center}
\[
\setlength{\arraycolsep}{3pt}
\begin{array}{r|rrlcrrr}
0&329^\circ&25'&48''.49''.50''.49''&=&329^\circ&25'&49''{,}00\\
1&25&10&39.46.38.40&=&25&10&40{,}75\\
2&80&55&40.35.33.29&=&80&55&34{,}25\\
5&248&10&21.10.16.16&=&248&10&15{,}75\\
9&111&9&58.53.47.46&=&111&9&51{,}00\\
19&308&38&48.47.52.52&=&308&38&49{,}75\\
27&34&38&7.\ 7.\ 4.\ 1&=&34&38&4{,}75\\
35&120&37&33.25.22.23&=&120&37&25{,}75\\
49&181&6&19.15.10.19&=&181&6&15{,}75\\
52&348&20&57.59.53.59&=&348&20&57{,}00
\end{array}
\quad
\GaussRightBrace{\begin{array}{rrr}
55^\circ&44'&51''{,}750\\&&53{,}500\\&&53{,}833\\&&53{,}812\\&&53{,}875\\
&&54{,}375\\&&55{,}125\\&&55{,}000\\&&53{,}750
\end{array}}\quad55^\circ44'\,54''{,}385.
\]
\begin{center}Falkenberg, heliotrope --- Hauselberg, stone.\end{center}
\[
\setlength{\arraycolsep}{3pt}
\GaussRightBrace{\begin{array}{r|rrlcrrr}
0&348&20&57.\ \text{etc.}&=&348&20&57{,}00\\
11&299&35&41.41.41.45&=&299&35&42{,}00
\end{array}}\quad28\ 17\ 42{,}275.
\]
\begin{center}Wilsede [heliotrope] --- Hauselberg, stone.\end{center}
\[
\setlength{\arraycolsep}{3pt}
\GaussRightBrace{\begin{array}{r|rrlcrrr}
0&299&35&41.\ \text{etc.}&=&299&35&42{,}00\\
6&134&52&30.29.21.25&=&134&52&26{,}25
\end{array}}\quad27\ 27\ 12{,}625.
\]
\medskip
\noindent\begin{minipage}[t]{.43\textwidth}
\begin{center}Simple angles.\end{center}
\[
\begin{array}{rrrlr}
122^\circ&36'&6''&\text{Hauselberg, stone}&\\
121&26&22&\text{Tempelbäume}&1\\
121&25&47&\text{---}&2\\
121&25&26&\text{---}&3\\
121&23&37&\text{---}&4\\
121&22&49&\text{---}&5\\
121&22&17&\text{---}&6\\
121&20&54&\text{---}&7
\end{array}
\]
\end{minipage}\hfill
\begin{minipage}[t]{.54\textwidth}
\vspace{2.7\baselineskip}
Several were not sighted; partly because it had already grown too late (the sun had long been below the horizon), partly because the Tempelbäume could not be distinguished clearly enough from the Wichelbäume, to which one or another of those already sighted might belong.
\end{minipage}

\clearpage
\noindent278\hfill\textsc{Correspondence.}\hfill\mbox{}\par\bigskip
\begin{center}
[6.]\\[7pt]
[On Station Adjustments.]\\[14pt]
\textsc{Gauss} to \textsc{Gerling}. Göttingen, 26 December 1823.
\end{center}
My letter reached the post too late and was brought back to me. I am therefore breaking it open again to add the practical instructions for elimination. Of course, there are various little advantages peculiar to particular cases that can be learned only ex usu.

I take your measurements on \textls[120]{Orber-Reisig} as an example[\textsuperscript{*})].

I first set
\[
\begin{array}{r@{\quad}r@{\;}c@{\;}l}
\text{[Direction to]}&1&=&0,
\end{array}
\]
\noindent next from $1\,.\,3$
\[
3=77^\circ57'\,53''{,}107
\]
\begin{center}(I prefer this because $1\,.\,3$ has more weight than $1\,.\,2$);\end{center}
\noindent then from
\[
\GaussRightBrace{\begin{array}{r|c|r@{\;}c@{\;}rrr}
13&1\,.\,2&2&=&26^\circ&44'&7''{,}423\\
50&2\,.\,3&2&=&&&6{,}507
\end{array}}\quad2=26^\circ44'\,6''{,}696;
\]
\noindent finally from
\[
\GaussRightBrace{\begin{array}{r|c|r@{\;}c@{\;}rrr}
26&1\,.\,4&4&=&136^\circ&21'&13''{,}481\\
6&2\,.\,4&4&=&&&8{,}529\\
78&3\,.\,4&4&=&&&11{,}268
\end{array}}\quad4=136^\circ21'\,11''{,}641.
\]
To refine the approximation a little further before proceeding, I find from

\vfill\noindent\rule{4cm}{.4pt}\par\smallskip
{\footnotesize
[\textsuperscript{*}) The angle measurements communicated by \textsc{Gerling} were (according to a sheet among \textsc{Gauss}' papers), where 1 denotes \textls[100]{Berger Warte}, 2 \textls[100]{Johannisberg}, 3 \textls[100]{Taufstein} and 4 \textls[100]{Milseburg}:
\[
\begin{array}{r@{\qquad}c@{\;}c@{\;}rrr}
\text{Rep.}&\multicolumn{5}{c}{\text{Angle}}\\
13&1\,.\,2&=&26^\circ&44'&7''{,}423\\
28&1\,.\,3&=&77&57&53{,}107\\
26&1\,.\,4&=&136&21&13{,}481\\
50&2\,.\,3&=&51&13&46{,}600\\
6&2\,.\,4&=&109&37&1{,}833\\
78\rlap{\hspace{.8em}\raisebox{.7ex}{\tiny.}}&3\,.\,4&=&58&23&18{,}161.\text{]}
\end{array}
\]
}

\clearpage
\noindent\mbox{}\hfill\textsc{Station Adjustments.}\hfill279\par\medskip
\[
\GaussRightBrace{\begin{array}{r|c|r@{\;}c@{\;}r}
13&1\,.\,2&1&=&-0''{,}727\\
28&1\,.\,3&1&=&0\\
26&1\,.\,4&1&=&-1{,}840
\end{array}}\quad1=-0''{,}855.
\]
Since any change common to all the directions is permissible as long as only their relative positions matter, I change all four by $+0''{,}855$ and set
\[
\begin{array}{r@{\;}c@{\;}rrr@{\;}l}
1&=&0^\circ&0'&0''{,}000&+a\\
2&=&26&44&7{,}551&+b\\
3&=&77&57&53{,}962&+c\\
4&=&136&21&12{,}496&+d.
\end{array}
\]
In the indirect method it is very advantageous to assign a correction to \textls[120]{each} direction. You can easily see this for yourself by working through the same example without this device, in which case you also lose the great convenience of always having the sum of the constant terms $=0$ as a check. I now form the four condition equations according to this scheme (when doing this myself and there are more terms, I generally separate the positive and negative terms), [where the constants are expressed in units of the third decimal place:]
\[
\begin{array}{rr@{\quad}rr@{\quad}rr@{\quad}rr}
ab&-1664&ba&+1664&ca&+23940&da&-25610\\
ac&-23940&bc&+9450&cb&-9450&db&+18672\\
ad&+25610&bd&-18672&cd&-29094&dc&+29094.
\end{array}
\]
The condition equations are thus:
\[
\begin{array}{r@{\;}c@{\;}rrrrr}
0&=&+6&+67a&-13b&-28c&-26d\\
0&=&-7558&-13a&+69b&-50c&-6d\\
0&=&-14604&-28a&-50b&+156c&-78d\\
0&=&+22156&-26a&-6b&-78c&+110d;
\end{array}
\]
\begin{center}Sum $=0$.\end{center}
To eliminate indirectly, I now note that if 3 of the quantities $a$, $b$, $c$, $d$ are set equal to 0, the fourth takes its largest value when $d$ is chosen as that quantity. Each quantity must of course be determined from its own equation, thus $d$ from the fourth. I therefore set $d=-201$

\clearpage
\noindent280\hfill\textsc{Correspondence.}\hfill\mbox{}\par\medskip
\noindent and substitute this value. The constant terms then become: $+5232$, $-6352$, $+1074$, $+46$; the rest stays the same.

I now take $b$ in turn, find $b=+92$, substitute and find the constant terms: $+4036$, $-4$, $-3526$, $-506$. I continue in this way until nothing remains to be corrected. In practice, however, I write down only the following scheme for this whole calculation:
\[
\setlength{\arraycolsep}{3pt}
\begin{array}{r|r|r|r|r|r|r|r}
&d=-201&b=+92&a=-60&c=+12&a=+5&b=-2&a=-1\\ \hline
+6&+5232&+4036&+16&-320&+15&+41&-26\\
-7558&-6352&-4&+776&+176&+111&-27&-14\\
-14604&+1074&-3526&-1846&+26&-114&-14&+14\\
+22156&+46&-506&+1054&+118&-12&0&+26.
\end{array}
\]
Since I carry the calculation only to the nearest $2000^{\mathrm{th}}$ [of an] arcsecond, I see that nothing now remains to be corrected. I therefore collect
\[
\begin{array}{r@{\;}c@{\;}r@{\qquad}r@{\;}c@{\;}r@{\qquad}r@{\;}c@{\;}r@{\qquad}r@{\;}c@{\;}r}
a&=&-60&b&=&+92&c&=&+12&d&=&-201\\
&&+5&&&-2&&&&&&\\
&&-1&&&&&&&&&\\ \cline{3-3}\cline{6-6}\cline{9-9}\cline{12-12}
&&-56&&&+90&&&+12&&&-201
\end{array}
\]
and add the Correctio communis $+56$, which gives:
\[
a=0\qquad b=+146\qquad c=+68\qquad d=-145,
\]
hence the values [of the directions]
\[
\begin{array}{r|rrr}
1&0^\circ&0'&0''{,}000\\
2&26&44&7{,}697\\
3&77&57&54{,}030\\
4&136&21&12{,}351.
\end{array}
\]
Almost every evening I make a new edition of the table, which is always easy to adjust further. Amid the monotony of measuring this always provides a pleasant diversion; one then also sees right away whether anything doubtful has crept in, what is still to be desired, etc. I recommend that you adopt this method. You will hardly ever eliminate directly again, at least when more than 2 unknowns are

\clearpage
\noindent\mbox{}\hfill\textsc{Station Adjustments.}\hfill281\par\medskip
\noindent involved. The indirect method can be carried out half asleep, or one can think about other things while doing it.

\noindent\hspace*{1.5em}.\ .\ .\ .\ .\ .

\vspace{2cm}
\begin{center}
\textsc{Gauss} to \textsc{Schumacher}. Göttingen, 22 December 1827.
\end{center}
The unit in my list of coordinates is $443{,}307885$ [Paris] lines; the logarithm for conversion to toises
\[
\llap{\raisebox{1.2ex}{\text{\tiny.}}\hspace{4cm}}=9{,}710\,1917.
\]
As things stand, the absolute scale rests solely on your baseline, or rather on the distance from Hamburg to Hohenhorn given to me by \textsc{Caroc}, $\log=4{,}141\,1930$, for which I have therefore taken: $4{,}431\,0013$. If, after the definitive calibration of your rods, your baseline and with it the distance given above from Hamburg to Hohenhorn should change, then all my coordinates will also have to be changed in the same ratio.

An important point in the way the work is organized is that a table is drawn up for each observation site, in which all the azimuths (in my sense of the term) are entered in order. Thus everything one needs from the observations is ready for the most convenient use, so that one refers back to the original records only in exceptional cases to settle this or that doubt. \ .\ .\ .\ .\ If the station differs from the target point, I do not in any way reduce the observations to the latter (centering), since they can be used just as conveniently without this reduction (that is, in the case of numerous intersections of subsidiary points that are themselves not stations).

\noindent\hfill\raisebox{.3ex}{\textbullet}\hspace*{1.2cm}\par
Drawing up such a table in turn depends on several considerations, and instructions for it can only be spread over several letters, so you might do well first to study this table yourself, as it were, and compare it with the observations, so that you can tell me more particu-
\vfill
\noindent{\small IX.}\hfill36
\endgroup
% End printed 281; the word particu- continues on printed 282.

